\section*{摘\quad 要}
\addcontentsline{toc}{section}{摘要}

携带卫星遥感先验底图的自主水下航行器（AUV）在现场巡测时，受机载算力、电池容量与作业时间的制约。由于卫星底图存在数日滞后与小尺度平滑偏差，航行器需在每个控制周期权衡：是开启高功耗图形处理器（GPU）在内部地图上深想，还是靠低功耗规划继续走并沿途采样，或是停下随水漂流原位复测。固定分配规则难以适应不同机载算力，单步元推理则会因规划器在有限样本上的过拟合而高估计算价值。本文提出一种基于留出想象海洋的集合同化与出发前预演方法。该方法利用八年东海卫星数据标定底图误差，围绕旧图生成48个想象中的海洋并随现场采样实时同化更新；其中32个用于规划航线，16个留出用于评判。出发前，母船或岸上计算机以4个留出海洋为假想真值，按机载算力与电池功耗模拟五种用时策略的24小时任务，选出平均遗憾最低的策略。在2018至2026年东海3197天卫星叶绿素数据构建的422个赤潮任务窗口上的实验表明：留出海洋在64至200,000次想象推演下与真实航线价值之差均不超过0.005；深想仅在基础算力过慢时有益，停下复测始终不划算；在148个测试任务上，出发前预演在慢计算机与中等计算机上分别取得0.276与0.265的时间平均遗憾，优于单步元推理，且与按开发年份调好的固定策略相差在置信区间内。

\vspace{1em}
\noindent\textbf{关键词：}海洋大数据；自主水下航行器；元推理；数据同化；赤潮巡测；计算与行动权衡

\clearpage
\section*{Abstract}
\addcontentsline{toc}{section}{Abstract}

Autonomous underwater vehicles (AUVs) carrying satellite prior maps must trade off computation against movement under onboard processor, battery, and mission-time constraints. Because satellite maps are several days old and smooth out small-scale features, the vehicle must decide at every control interval whether to activate a high-power graphics processing unit (GPU) to plan deeply, continue cruising and sampling under low-power background planning, or stop propulsion to drift and resample locally. Fixed time-allocation rules do not generalise across processor speeds, while classic one-step metareasoning overestimates the value of computation by trusting the planner's own search trace on finite ensemble members. We propose an ensemble data assimilation and pre-mission rehearsal method using held-out imagined oceans. Using eight years of satellite chlorophyll-a data over the East China Sea to calibrate stale-map uncertainty, the vehicle generates 48 imagined oceans around the stale map and updates them with in-situ fluorometer samples via a localised ensemble adjustment Kalman filter. Thirty-two members are searched by a cross-entropy route planner and 16 are held out for unbiased evaluation. Before launch, a shipboard or shore computer treats 4 held-out members as synthetic truths and simulates complete 24-hour missions for five time-allocation strategies under the vehicle's actual compute speed and battery budget, selecting the strategy with the lowest mean regret. Across 422 spring bloom windows from 3,197 daily maps (2018--2026), held-out evaluation stays within 0.005 of true route value from 64 to 200,000 rollouts, GPU planning helps only when baseline compute is too slow, and stopping never pays. On 148 test missions, rehearsal achieves time-averaged regrets of 0.276 and 0.265 on slow and moderate computers, outperforming one-step metareasoning and matching development-tuned strategies within confidence intervals without prior missions on that vehicle.

\vspace{1em}
\noindent\textbf{Keywords: }Ocean big data; Autonomous underwater vehicle; Metareasoning; Data assimilation; Harmful algal bloom survey; Thinking-acting trade-off
